\section{Evaluation}
\label{sec:eval}
The primary campaign uses the official 1.8.2 ARM64 image and the implemented
2.0 release-intent semantics. Historical competitor measurements remain
separate. Our evaluation asks whether planning matches a reference domain,
whether recorded progress survives injected failures, and what planning costs
on larger graphs. A supplemental stage-order experiment addresses the running
example directly; its setup and results are reported below.

\subsection{Study Design and Evidence Boundaries}
\label{sec:eval:diff}
The reference oracle is approximately 300 lines of Python. It predates the
complete formal model and implements a restricted fragment. A seeded generator
creates real Git repositories with 5--20 packages, layered DAGs, 4--12 intent
units, and mixed release-tag states. Golden cases gate the comparison between
the oracle and the released binary. Agreement tests that shared fragment;
it does not prove the implementation refines every definition in
Section~\ref{sec:model}.

Two profiles answer different questions. The core profile stays within the
shared semantic domain. The full profile includes prerelease baselines,
channel-only entries, and directives outside the oracle's complete coverage.
Full-profile differences are diagnostic observations, not passing conformance
checks or automatically established implementation bugs. We retain rejected
plans, their diagnostics, and the unit of counting: a history can have several
differing package rows.

Recovery campaigns execute interrupted releases and inspect subsequent plans,
versions, and tags. Separate fixtures introduce refused publication and a
competing Git push. These observations test selected implementation traces,
not every interleaving of external writes. The timing campaign varies graph
size, shape, intent-unit count, and unrelated history. It uses a contended
host capped at two CPUs and 4 GiB; it is not a controlled cross-tool benchmark.
All reported primary counts were reaggregated from saved records for this
revision. That integrity check is not a new execution of the original campaign.

\subsection{Planner Agreement and Specification Regressions}
\label{sec:eval:fresh}
\begin{table}[t]
\centering
\caption{Two differential profiles. Diagnostic equality is not executable
conformance: four matching full-profile cases were refused.}
\label{tab:oracle}
\small
\begin{tabular}{@{}lrr@{}}
\toprule
Observation & Core & Full \\
\midrule
Generated histories & 1,500 & 1,500 \\
Accepted plans & 1,500 & 997 \\
E195 diagnostic refusals & 0 & 503 \\
Matching comparisons & 1,500 & 550 \\
Differing package rows & 0 & 1,914 \\
\bottomrule
\end{tabular}
\end{table}

All 1,500 core instances agreed, with no invalid sample. Of 1,500 full-profile
histories, 997 produced accepted plans and 503 produced complete E195
refusals for proposed regressive versions (Table~\ref{tab:oracle}). Among
accepted plans, 546 comparisons were clean. Four refused dossiers also had no
row difference, giving 550 matching diagnostic comparisons overall. Each
proposed regression set in the refused dossiers was independently checked by
the campaign's validation logic. Rejected diagnostic output is not an
executable release plan.

The 1,914 differing rows retain the campaign's heuristic classes: PG1 862,
PG2 375, PG2x 14, and PG3? 663, with UNKNOWN 0. These labels partition
observations for inspection; UNKNOWN 0 does not show that the oracle explains
every current semantic rule correctly. The full-profile results therefore
limit the refinement claim rather than extend the 1,500 core agreements to
all prerelease behavior.

The binary passed six selected specification vectors, 133--138. They exercise
fresh discharge, suppression, channel eligibility and ordering, inherited
channel drift, and source exclusion. Vector 135 observes a downward W208
diagnostic, not the upward audit API. These targeted cases complement random
core instances but are not complete specification conformance coverage.

\subsection{Recovery and Historical Replay}
Two isolated orphan repetitions retained the failed consumer's version and
reached an empty third plan. Across two broader recovery campaigns, all
24 episodes converged within four runs, with zero structural G3/G5 check
failures. These checks concern the tested schedules and predicates. They do
not establish the liveness premise under permanent registry failure, nor do
they make an unrecorded external write unambiguous.

Three fresh 1.8.2 fault cells passed 8 orphan checks, 12 clean competing-push
checks, and 14 conflicting-push checks. The latter distinguish preservation
of completed release tags from incorporation of a concurrent source change.
Different scenarios have different assertions; their counts are not a common
quality score. The artifact retains the commands, logs, and verdicts needed
to see what each assertion establishes.

Five-release historical replay matched 36 of 42 versions in
\code{google-cloud-node} and 7 of 64 in \code{conventional-changelog}. Those
projects did not use the modeled intent language. Their layouts and release
policies were reconstructed for the replay, and mismatches may reflect that
reconstruction as well as different release conventions. Consequently the
experiment demonstrates limits of policy transfer; it is not evidence that
our protocol reproduces arbitrary tools' historical version decisions.

\subsection{Comparison with Pinned Release Configurations}
\label{sec:compare}
The historical refusal and competing-push fixtures pin Lerna 10.0.1, Nx
23.1.2, and Changesets 3.0.1. In those configurations, Lerna and Nx tagged the
denied version before its failed upload, while Changesets left it untagged
and recovered by querying npm. The observed propagation policies were
unconditional for Lerna/Nx, range-gated for Changesets, and per-commit for our
protocol. These are observations of the pinned configurations, not universal
properties of every mode those tools expose.

In the competing-push fixture, compared runs left tags outside branch
ancestry; our historical run merged before its final push and planned the
colleague's commit next. The new 1.8.2 fault cells above repeated our side,
not the competitors. Release-please's tag timing is characterized from its
documented workflow, not an executed comparison~\cite{releaseplease}.
No cross-tool speed or aggregate superiority claim follows from these tests.
The comparison isolates ledger and recovery boundaries, which differ from
ordinary DAG scheduling.

\subsection{Planning Cost}
\begin{table}[t]
\centering
\caption{Saved planning measurements. $P$: packages; $U$: intent units;
$C_0$: padding commits. Times are seconds; RSS is maximum MiB. Each cell
reports the median and observed minimum--maximum, not a confidence interval.}
\label{tab:scale}
\small
\setlength{\tabcolsep}{3.5pt}
\begin{tabular}{@{}rlrrrlr@{}}
\toprule
$P$ & Shape & $U$ & $C_0$ & Median & Range & RSS \\
\midrule
100 & layer & 200 & 0 & .017 & .017--.017 & 15 \\
1,000 & layer & 200 & 0 & .250 & .235--.252 & 21 \\
10,000 & layer & 200 & 0 & 19.272 & 18.992--20.000 & 92 \\
1,000 & chain & 200 & 0 & .254 & .241--.258 & 26 \\
10,000 & chain & 200 & 0 & 19.354 & 18.933--19.645 & 118 \\
1,000 & star & 200 & 0 & .244 & .237--.247 & 19 \\
1,000 & layer & 1,000 & 0 & .290 & .280--.302 & 26 \\
1,000 & layer & 200 & 5,000 & .664 & .640--.665 & 30 \\
\bottomrule
\end{tabular}
\end{table}

All eight cells are valid observations under the recorded host conditions
(Table~\ref{tab:scale}). Layered graphs have 185, 2,002, and 19,919 edges at
100, 1,000, and 10,000 packages; chain and star graphs have $P-1$ edges.
The 10,000-package costs are substantially above the 1,000-package cells.
We do not infer linear end-to-end scaling from the propagation bound:
repository construction, history access, scope resolution, and output work
are separate costs. Increasing intent to 1,000 units and adding 5,000 padding
commits probe different inputs and should not be merged into a package-count
trend. Timing is sensitive to the contended host and changed harnesses, so
these observations are not compared with earlier campaign latencies.

The fixtures passed \code{git fsck}. The optimized \code{commit-tree}
generator matched an ordinary small fixture's commit IDs, references,
messages, and log output. This check reduces a fixture-generation threat but
does not turn synthetic graphs into a representative sample of industrial
repositories. Their purpose is to expose costs and graph shapes explicitly.

\subsection{Counterexamples and Test Obligations}
\label{sec:eval:formal}
The formalization identifies six failures of simpler rules: reusing the train
window for discharge; changing suppression during retries; allowing inherited
channels to drift; implementing upward audits with reverse reachability alone;
resolving bumps before channels; and merging unit identities during shared
traversal. The motivating orphan is a further direct consequence of
source-based rather than consumer-based pending work.

These examples explain necessary distinctions in the model. They are not
measurements of defect prevalence and are not assertions of outstanding bugs
in the released specification. The targeted vectors exercise selected cases;
we do not equate the number of logical counterexamples with the number of
independently reproduced implementation defects. In particular, the inherited
channel example in Remark~\ref{rem:drift} shows why stable completed versions
alone do not establish retry-version stability.
\subsection{Mixed-Chain Availability, Recovery, and CI}
The fresh 1.9.0 Linux ARM64 campaign executes both the three-package
npm $\to$ Docker $\to$ Docker chain and the four-package running example.
Build outputs carry fixture content through OCI archives; private Verdaccio
and Docker registries distinguish local output from publication~\cite{oci,buildkit}.
Each topology repeats five configurations six times in seeded order: mixed and
serialized availability, a missing barrier, refused publication, and failed building.

\begin{table}[t]
\centering
\caption{Fresh 1.9.0 executions. Designed failures are retained controls;
checks are verifier predicates, not independent samples.}
\label{tab:mixed}
\small
\setlength{\tabcolsep}{2pt}
\begin{tabular}{@{}lrrl@{}}
\toprule
Scope & Runs & Checks & Initial outcomes \\
\midrule
Three-package chain & 30 & 522 & 12 pass; 18 fail \\
Four-package chain & 30 & 798 & 12 pass; 18 fail \\
CI entry point & 3 & 81 & 3 pass \\
\bottomrule
\end{tabular}
\end{table}

Per topology, 12 repairable failures completed with unchanged versions; 24
completed cases then reran without build or publication. Six missing-barrier
controls were not repaired. Enabled barriers placed app building after base
publication; without one, the app failed against the unavailable image. In all
six mixed three-package runs, base building overlapped npm publication.

In the four-package mixed runs, UI building overlapped delayed assets
publication. All six base builds finished before UI publication started. This
is a valid local-output schedule: base building need not await UI publication,
but overlap is not required. Intervals use a shared monotonic clock~\cite{pytime}.
An initial campaign attempt incorrectly required
strict overlap and failed its oracle. We corrected that predicate before a
full rerun from fresh state; a focused oracle test accepts completion before publication
starts and rejects a build starting after publication ends. The correction and
earlier failure are retained, not counted as implementation evidence.

The three CI repetitions kept the workflow entry point and shared configuration
unchanged while adding packages. They selected required tests through the graph,
blocked release stages for an injected or missing test, propagated a specification
change to its consumer, and observed a final no-op release rerun. These fixtures
support selected ordering, recovery, and CI-integration paths on one Linux ARM64
environment, not a universal schedule, reliability rate, throughput result, or
independent adoption study.

